% Independent replacement proof supplied by /root/homotopy_levels, 2026-11-19.
% Replacement body for C4.05; retain its statement identifier and heading.
Use the actual finite proper conical deployment with compatible iterated fibre collars of C2.05, in the metric setting of C4.01. We spell out the radial data used in this assertion. Each stratum \(S_i\) has a fixed control tube \(N_i\), a retraction \(\pi_i:N_i\to S_i\), and a radius \(\rho_i:N_i\to[0,1)\) with zero set \(S_i\). Each closed subtube \(0\le\rho_i\le a<1\) is proper over \(S_i\times[0,a]\) under \((\pi_i,\rho_i)\). This follows from the proper compact link bundle in its mapping-cone presentation: the product of that bundle with \([0,a]\), followed by the cone quotient, is proper.

On the resolution \(\beta_j:\Sigma_j\to\overline{S_j}\), the interior is identified with \(S_j\). If the original construction retains a cut core instead, its collar stretching gives the same identification: append an outer collar with coordinate in \((-1,0]\), combine it with an inner collar \([0,a)\), and choose a monotone homeomorphism \((-1,a)\to(0,a)\) equal to the identity near \(a\). Extend by the identity outside the collar. At corners apply these reparametrizations in the commuting product coordinates. Retained outer faces in intermediate induction diagrams are handled by C1.B2a. The compatible iterated fibre collars give the following precise simultaneous coordinates. For
\(d=(i_0<\cdots<i_{n-1}<j)\), put
\[
 F_d=\bigcap_{k<n}\Sigma_{i_kj}=\Sigma_d,
 \qquad F_d^\circ=F_d\setminus
       \bigcup_{\ell\notin\{i_0,\ldots,i_{n-1},j\}}
                  (F_d\cap\Sigma_{\ell j}).
\]
Thus all remaining incoming boundary faces, including intermediate ones not selected by \(d\), are removed in \(F_d^\circ\). After fixing compatible collar widths, the intersection of the selected collar neighborhoods in the interior is identified with
\[
 c_d:F_d^\circ\times(0,a)^n\xrightarrow{\cong}
           S_j\cap\bigcap_{k<n}N_{i_k}.
\]
Here and below \(c_d\) includes the blowdown on the interior. The widths may first be normalized by the given collar coordinates. Their compatibility with the boundary fibrations gives continuous maps
\[
 p_k:F_d^\circ\times(0,a)^k\longrightarrow S_{i_k}
\]
with
\[
 \rho_{i_k}c_d(z,r)=r_k,
 \qquad
 \pi_{i_k}c_d(z,r)=p_k(z,r_0,\ldots,r_{k-1}).
 \tag{C4.05a}
\]
The absence of later radii in the second formula is essential: the later collars are vertical for the earlier source fibration. These are the simultaneous product and fibre-compatibility axioms of the actual resolved collar system, equivalently the compatible radial control constructed in C2.05. Pairwise commutativity of arbitrary retractions alone would not supply (C4.05a).

Here is the translation from the resolved data of C2.05. On the face \(\Sigma_{i_kj}\), the faces with labels \(i_\ell<i_k\) are the inverse images of the corresponding boundary faces of the base \(\Sigma_{i_k}\); this is the flag inverse-image identity. Their collars are carried to the base collars. The faces with labels \(i_\ell>i_k\) are vertical faces of the source fibre, and their compatible collars leave that source map fixed. Applying the commuting collar maps successively therefore carries the source only through the earlier coordinates, exactly as in (C4.05a). The transverse product-collar axiom identifies intersections of the chosen collar neighborhoods with the simultaneous collar of their face intersection, after shrinking the finite collar system. It also preserves every nonselected face. Removing all boundary faces to pass to \(S_j\) consequently removes the nonselected faces from the base \(F_d\), giving precisely \(F_d^\circ\), and requires each selected radius to be positive. This proves the displayed coordinate description from the full iterated fibre-collar data, rather than assuming it for arbitrary retractions.

\textbf{A cofinal family consisting of actual neighborhoods.} For every positive continuous function \(\epsilon_i:S_i\to(0,a)\), let
\[
 T_i(\epsilon_i)=
 \{x\in N_i:\rho_i(x)<\epsilon_i(\pi_i(x))\}.
\]
These are genuine open neighborhoods of \(S_i\) in the required open complement of its lower frontier. They form a cofinal family. Indeed, for a prescribed neighborhood \(U_i\), choose a fixed smaller closed subtube \(K_i=\{\rho_i\le a'\}\), \(0<a'<a\). The set \(K_i\setminus U_i\) is closed; properness makes its image under \((\pi_i,\rho_i)\) a closed subset of \(S_i\times[0,a']\) disjoint from the zero section. There is a positive continuous function \(\epsilon_i<a'\) whose subgraph misses that image. To construct it explicitly, choose a cover of \(S_i\) by open sets \(B_\alpha\) and numbers \(b_\alpha>0\) for which \(B_\alpha\times[0,b_\alpha]\) misses the closed image. A locally finite subordinate partition of unity gives \(\epsilon_i(y)=\frac12\sum_\alpha\lambda_\alpha(y)b_\alpha\), capped at \(a'/2\). Each active bound is admissible at \(y\), hence so is their weighted average. Then \(T_i(\epsilon_i)\subset U_i\).

Consequently the germ \(V_d\cap S_j\) is represented by the actual intersections
\[
 H_\epsilon=S_j\cap\bigcap_{k<n}T_{i_k}(\epsilon_{i_k}).
\]
The last neighborhood, that of \(S_j\) itself, imposes no restriction after intersecting with \(S_j\). In the fixed simultaneous coordinates, (C4.05a) identifies \(H_\epsilon\) with
\[
 \left\{(z,r):
 0<r_k<\epsilon_{i_k}\bigl(p_k(z,r_0,\ldots,r_{k-1})\bigr)
 \quad(0\le k<n)\right\}.\tag{C4.05b}
\]
These are recursively varying radius regions. The family is cofinal in both required senses: every member is exactly an actual intersection of the indicated neighborhood representatives, and the preceding properness argument puts some such member inside every prescribed tuple of neighborhoods. We do \emph{not} claim that arbitrary positive functions on the entire resolved face or on fibre directions produce cofinal ambient neighborhood systems.

\textbf{The homotopy equivalence and its compatibility.} Each region (C4.05b) is homeomorphic to \(F_d^\circ\times(0,1)^n\). The homeomorphism is given successively by
\[
 r_k=\epsilon_{i_k}\bigl(p_k(z,r_0,\ldots,r_{k-1})\bigr)u_k,
 \qquad 0<u_k<1.
\]
Its inverse divides by the same strictly positive continuous functions, successively in increasing \(k\). Projection to \(z\) is therefore a homotopy equivalence. All the transition inclusions for smaller control functions commute with this one fixed projection. We obtain a natural objectwise equivalence from the cofinal diagram of the \(H_\epsilon\) to the constant space \(F_d^\circ\), and hence an equivalence of prospaces, with full coherent compatibility. There is no inference merely from unrelated isomorphisms in \(\mathrm{HoTop}\).

Finally \(F_d^\circ\hookrightarrow F_d\) is a homotopy equivalence. The finitely many remaining commuting collars can be pushed inward, using in each coordinate a function positive at zero and equal to that coordinate outside a smaller collar. They preserve the selected faces because their collars commute. The resulting global homotopy has image in \(F_d^\circ\) at its final time and preserves \(F_d^\circ\) at every time. Thus
\[
 V_d\cap S_j\simeq F_d^\circ\simeq F_d=\Sigma_d.
\]
C4.01 supplies \(V_d\simeq V_d\cap S_j\). Empty face intersections have an empty cofinal neighborhood intersection; the same statement then holds with the empty space. For a one-element flag the assertion reduces to the inward-collar equivalence between the resolved piece and its stratum.

The rectangular version must be qualified. An arbitrary neighborhood of an intermediate stratum can impose \(r_1<r_0^2\) near a lower stratum, so no product rectangle with two fixed positive radii is contained in it. Formula (C4.05b), with radius functions pulled back from the actual stratum bases, is what makes the comparison valid. The resulting equivalence is canonical in the homotopy category for the fixed compatible deployed collar germs; the contractible choices of smaller radius functions and the displayed natural projection specify its compatible homotopies.
